\chapter{Konstruktionen totaler Ordnungen}
\label{b15guide:ch:constructions}

\section{Einschränkungen und Umkehrung}

\begin{theorem}[Erhaltung der vollständigen Vergleichbarkeit]
\label{b15guide:thm:restriction-dual}
Für jede Teilmenge $S\subseteq A$ ist die eingeschränkte
Relation $\leq_S\ =\ \leq\cap(S\times S)$ eine totale
Ordnung auf $S$.
Auch die umgekehrte Relation
\[
 x\leq^{\mathrm{op}}y\quad\Longleftrightarrow\quad y\leq x
\]
ist eine totale Ordnung auf $A$.
\end{theorem}

Die Begriffe Restriktion und Umkehrung stammen aus Band~12.
Hier wird nur ihre Erhaltung der Totalität festgehalten.

\section{Geordnete Summen und lexikographische Produkte}

\begin{definition}[Geordnete Summe]
\label{b15guide:def:sum}
Für totale Ordnungen $(A,\leq_A)$ und $(B,\leq_B)$
verwenden wir die disjunkte Vereinigung
\[
 A+B=(\{0\}\times A)\cup(\{1\}\times B).
\]
Innerhalb jedes Teils gilt dessen ursprüngliche Ordnung.
Jedes $(0,a)$ wird vor jedes $(1,b)$ gesetzt.
Dies definiert die Relation $\leq_+$.
\end{definition}

\begin{theorem}[Totalität der geordneten Summe]
\label{b15guide:thm:sum}
$\leq_+$ ist eine totale Ordnung auf $A+B$.
Der erste Teil ist ein Anfangsabschnitt, der zweite ein
Endabschnitt. Ist einer der Träger leer, besteht die Summe nur aus
dem markierten anderen Teil mit dessen ursprünglichen Vergleichen.
\end{theorem}

\begin{definition}[Lexikographische Ordnung]
\label{b15guide:def:lex}
Für totale Ordnungen $(A,\leq_A)$ und $(B,\leq_B)$
setzen wir auf $A\times B$
\[
 (a,b)\leq_{\mathrm{lex}}(a',b')
 \quad\Longleftrightarrow\quad
 \bigl(a<_A a'\ \lor\ (a=a'\land b\leq_B b')\bigr).
\]
Die erste Koordinate hat Vorrang.
\end{definition}

\begin{theorem}[Totalität des lexikographischen Produkts]
\label{b15guide:thm:lex}
$\leq_{\mathrm{lex}}$ ist eine totale Ordnung auf $A\times B$.
Für festes $a\in A$ ist $\{a\}\times B$ ordnungskonvex;
für $a<_A a'$ liegt jedes Element von $\{a\}\times B$
vor jedem Element von $\{a'\}\times B$.
\end{theorem}

\begin{definition}[Geordnete Summe einer Familie]
\label{b15guide:def:indexed-sum}
Sei $(I,\leq_I)$ total geordnet, und sei für jedes $i\in I$
eine totale Ordnung $(A_i,\leq_i)$ gegeben. Auf
\[
 \sum_{i\in I}A_i=\{(i,a)\mid i\in I,\ a\in A_i\}
\]
definieren wir
\[
 (i,a)\preceq(j,b)
 \quad\Longleftrightarrow\quad
 \bigl(i<_I j\ \lor\ (i=j\land a\leq_i b)\bigr).
\]
Die Ordnung $\leq_i$ wird im zweiten Fall, also innerhalb
desselben Teils, ausgewertet.
\end{definition}

\begin{theorem}[Totalität der geordneten Familiensumme]
\label{b15guide:thm:indexed-sum}
$\preceq$ ist eine totale Ordnung auf $\sum_{i\in I}A_i$.
Jeder Teil $\{i\}\times A_i$ ist ordnungskonvex;
die verschiedenen nichtleeren Teile stehen in der Reihenfolge
ihrer Indizes. Leere $I$ und leere $A_i$ sind erlaubt.
Die Ordnungen der Teile und des Indexträgers werden als gegeben
vorausgesetzt.
\end{theorem}
